\section{Dimension and trace fluctuations}\label{spm:sec:fluctuations}
Choose a matrix uniformly from $\mathcal A_n$, and write $D_n,T_n$ for its dimension and diagonal sum. Their joint moment generating function is
\[
 \E e^{aD_n+bT_n}=\frac{A_n(e^a,e^b)}{A_n(1,1)}.
\]
Theorem~\ref{spm:thm:main} applies in a fixed complex neighborhood of $(a,b)=(0,0)$, including derivative remainders. It therefore supplies all fixed cumulants by differentiation.

\subsection{The involution factor}
The one-dimensional saddle of Lemma~\ref{spm:lem:involution}, with Stirling's formula for $n!$, gives uniformly near $v=1$
\begin{align}\label{spm:eq:logI}
 \log I_n(v)={}&\frac n2(\log n-1)+v\sqrt n-\frac{v^2}{4}
              -\frac12\log2\notag\\
 &+\left(\frac{v^3}{24}+\frac v4\right)n^{-1/2}
   -\left(\frac{v^2}{16}+\frac1{12}\right)n^{-1}
   +O(n^{-3/2}).
\end{align}
For completeness, expand $\rho(h)$ from \eqref{spm:eq:scaling} in
$vr+r^2/2-n\log r$, expand the angular normalizer using $H$ in \eqref{spm:eq:H}, and add
\[
 \log n!=(n+\tfrac12)\log n-n+\tfrac12\log(2\pi)
             +\frac1{12n}+O(n^{-3}).
\]
At each order this uses finitely many even Gaussian moments. The same Taylor bound as \eqref{spm:eq:taylorbound}, with no gamma variable, proves the remainder. Thus \eqref{spm:eq:logI} is also an all-fixed-order construction. At $v=1$ its logarithmic coefficient of $n^{-1}$ is $-7/48$.

\begin{proposition}[First joint cumulants]\label{spm:prop:cumulants}
With $L=\log2$,
\begin{align}
 \E D_n={}&\frac{n}{2L}+\frac1{2L}-\frac12-\frac L4+O(n^{-1/2}),\label{spm:eq:ED}\\
 \Var D_n={}&\frac{n(1-L)}{4L^2}
 +\frac{1-L}{4L^2}+\frac{L-1}{8}+O(n^{-1/2}),\label{spm:eq:VD}\\
 \E T_n={}&\sqrt n+L-\frac12
 +\frac{3/8-L+L^2/2}{\sqrt n}+O(n^{-1}),\label{spm:eq:ET}\\
 \Var T_n={}&\sqrt n+2L-1
 +\frac{5/8-4L+5L^2/2}{\sqrt n}+O(n^{-1}),\label{spm:eq:VT}\\
 \Cov(D_n,T_n)={}&-\frac12+\frac{1-L}{2\sqrt n}+O(n^{-1}).\label{spm:eq:cov}
\end{align}
\end{proposition}
\begin{proof}
Take the logarithm of \eqref{spm:eq:main}, using
$C_1n^{-1/2}+D_2n^{-1}$ for its first normalized logarithmic terms, and insert \eqref{spm:eq:logI}. Differentiate with respect to $a$ and $b$. The required identities are
\[
 \left.\frac{\dd L(e^a)}{\dd a}\right|_{a=0}=-\frac12,
 \qquad
 \left.\frac{\dd^2 L(e^a)}{\dd a^2}\right|_{a=0}=\frac14,
 \qquad \partial_b=v\partial_v.
\]
Both $-(n+1)\log L(e^a)$ and $-\log(1+e^a)$ must be retained. The uniform holomorphic remainder justifies the differentiated error bounds.
\end{proof}

\begin{theorem}[Two independent Gaussian scales]\label{spm:thm:gaussian}
Let $\sigma_D^2=(1-L)/(4L^2)>0$. Then
\begin{equation}\label{spm:eq:gaussian}
 \left(\frac{D_n-n/(2L)}{\sigma_D\sqrt n},
       \frac{T_n-\sqrt n}{n^{1/4}}\right)
 \ \Longrightarrow\ (Z_1,Z_2),
\end{equation}
where $Z_1,Z_2$ are independent standard normal variables.
\end{theorem}
\begin{proof}
Substitute $a=s/(\sigma_D\sqrt n)$ and $b=t/n^{1/4}$ in the logarithm of the marked ratio. After centering, its limit is $(s^2+t^2)/2$ on a neighborhood of the origin. The leading $a$ term comes from $-(n+1)\log L(e^a)$, the leading $b$ term from $e^b\sqrt n$, and all mixed terms vanish on these scales. Convergence of the joint moment generating function proves the result.
\end{proof}

The nonzero limiting \emph{unscaled} covariance in \eqref{spm:eq:cov} is compatible with independent scaled limits. There is also an exact lattice constraint:
\begin{equation}\label{spm:eq:parity}
 T_n\equiv n\pmod2.
\end{equation}
No span-one local limit or total-variation approximation of $T_n$ by an unrestricted Poisson law is asserted. The valid local analytic residue is instead
\begin{equation}\label{spm:eq:trace-residue}
 \E[v^{T_n}]e^{-(v-1)\sqrt n}
 \longrightarrow \exp\left(\frac{(2L-1)(v^2-1)}4\right)
\end{equation}
uniformly near $v=1$. This follows by combining the $v$-dependent terms in \eqref{spm:eq:logI} with $S(L,v)-S(L,1)$.
